\section{Outstanding tasks}
\label{sec:outstanding}

The established endpoint is the complete calibrated response theorem,
together with its domains, finite-order estimates, exact zero
preservation, and conjugated dynamics. The next work concerns what can
be recovered from that response and whether the canonical source
supports a further estimate.

\paragraph{Sharp finite-window first-jet conditioning.}
The decision recorded as TND-A-B02 \cite{TNDAuxConditioning} and the
frozen TND-A-E03 request \cite{TNDAuxE03Request} ask for the exact
worst-case recovery error for \eqref{eq:first-jet}. The parameters
$M,K,T$ are finite, $s>0$ is fixed, and the norms and error units must
be stated before comparing proper times. For an additive Euclidean
error ball on the observed pair, the question is to determine the
sharp pointwise and window bounds governed by
\begin{equation}
 \norm{(mP_a)^{-1}}_2
 =q(\tau)^{-1}\ee^{\tau k^2}\norm{P_a^{-1}}_2,
 \qquad |k|\le K,\quad 0\le\tau\le T.
 \label{eq:pending-conditioning}
\end{equation}
Sharpness over unrestricted pair errors must be distinguished from
sharpness over perturbations realizable by admissible theta sources
or globally consistent differentiated error functions. The same
analysis must track the transformed error ellipsoid in
\eqref{eq:covariant-pair} and compare it with uncertainty propagated
from the same source. The present manuscript supplies the exact map;
the requested sharp theorem has not been executed or assessed here.

\paragraph{Assessment of the conditioning result.}
The planned TND-A-E04 assessment should check the completed E03 result
against its fixed error model, endpoints, and sharpness class. Neither
a convenient coordinate choice nor a smaller plotted amplitude is
evidence of improved recovery. This diagnostic has no assumed
physical noise law and no built-in positive collision margin.

\paragraph{A source-specific consequence.}
A consequential next estimate would have to use the exact canonical
theta structure quantitatively. In particular, the present positivity,
regularity, reflection, and heat identities do not by themselves bound
the simultaneous smallness of $H_\lambda$ and $H_\lambda'$ near a
potential collision. A proposed lower bound or signed inequality must
state its parameter range and identify precisely which premise fails
for the relevant colliding controls. If the conditioning diagnostic
shows only attenuation and reversal, that would justify stopping this
representation as an estimating route unless an additional
source-specific mechanism is found.

\paragraph{Intrinsic preparation and physical implementation.}
The scan $\beta(u)=\pi\ee^{4u}$, processing operator
$(\partial_u^2-1)/8$, spatial imprint, linear amplitude gain, matched
passive bank, and phase-sensitive detector law are supplied. Deriving
any of them from a specified physical apparatus would be a separate
task with additional assumptions. Such a derivation should also
distinguish Euclidean evolution from real unitary laboratory time and
account for preparation, calibration, and detection errors.

\paragraph{Independent mathematical review and supporting evidence.}
The current proof notes and their assessments are internal LLM-assisted
work with shared premises. Independent specialist review should
examine the geometric conventions, complete-domain construction,
weighted zero-time argument, and absolute calibration before stronger
reliance is placed on the representation. Numerical checks may support
finite-window constants and truncation bounds after the analytic
questions are fixed. They would remain supporting evidence: a finite
computation cannot replace a uniform estimate or a collision exclusion.
